\documentclass[a4paper,11pt]{article}
\usepackage[utf8]{inputenc}
\usepackage[T1]{fontenc}
\usepackage[french]{babel}
\usepackage[left=2cm,right=2cm,top=2cm,bottom=2cm]{geometry}
\usepackage{xcolor}
\usepackage{hyperref}
\usepackage{fontawesome5}
\usepackage{parskip}

\definecolor{primary}{RGB}{0, 140, 80}
\hypersetup{colorlinks=true, urlcolor=primary}

\begin{document}

\noindent
\begin{minipage}[t]{0.5\textwidth}
    \textbf{Lo\"ic Aron MBASSI EWOLO}\\
    \'El\`eve-Ing\'enieur Bac+5 -- Double dipl\^ome ENSEA\\[3pt]
    \faMapMarker\ Cergy, France\\
    \faPhone\ (+33) 7 59 52 33 98\\
    \faEnvelope\ \href{mailto:loicaron.mbassiewolo@ensea.fr}{loicaron.mbassiewolo@ensea.fr}
\end{minipage}
\hfill
\begin{minipage}[t]{0.4\textwidth}
    \raggedleft
    Cergy, le 21 mars 2026
\end{minipage}

\vspace{1.2cm}

\hfill
\begin{minipage}[t]{0.5\textwidth}
    \raggedleft
    \textbf{ENGIE}\\
    Service Recrutement\\
    La Garenne-Colombes, France
\end{minipage}

\vspace{1cm}

\noindent\textbf{Objet :} Candidature en alternance -- Data Scientist -- Septembre 2026

\vspace{0.8cm}

Madame, Monsieur,

La pr\'ediction de la consommation \'energ\'etique \`a fine granularit\'e -- par b\^atiment, par heure, selon les conditions m\'et\'eo et les comportements d'usage -- est un probl\`eme ML dont la pr\'ecision se traduit directement en optimisation des flux d'\'electricit\'e et en r\'eduction des co\^uts d'approvisionnement. ENGIE, en structurant des \'equipes data science d\'edi\'ees \`a la transition \'energ\'etique, offre un contexte o\`u la qualit\'e des mod\`eles a un impact environnemental mesurable.

Mon stage \`a l'INRIA Saclay (\'equipe TRIBE, avril-ao\^ut 2026) me forme \`a la rigueur exp\'erimentale n\'ecessaire : fine-tuning de mod\`eles, analyse de l'impact des hyperparamètres, \'evaluation robuste et suivi MLflow. Mon exp\'erience au MILA (\'et\'e 2025) sur la dynamique de convergence des r\'eseaux profonds me donne les outils th\'eoriques pour diagnostiquer les \'echecs des mod\`eles en production. Mon projet Urban Waste Management (1\textsuperscript{er} prix CITS) illustre ma capacit\'e \`a d\'eployer des mod\`eles ML sur donn\'ees IoT avec monitoring des s\'eries temporelles en temps r\'eel.

Ma formation en math\'ematiques appliqu\'ees (Polytechnique Yaound\'e) -- statistiques bay\'esiennes, probabilit\'es, optimisation -- fournit les fondements analytiques pour aborder les probl\`emes de r\'egression temporelle et de d\'etection d'anomalies sur les r\'eseaux \'electriques. Mon UROP 2024 (Google DeepMind) d\'emontre ma capacit\'e \`a r\'eduire significativement les temps de calcul par fine-tuning syst\'ematique.

Je suis disponible d\`es septembre 2026 pour cette alternance \`a La Garenne-Colombes. Travailler sur des probl\'ematiques data science \`a fort enjeu environnemental chez ENGIE est une perspective concr\`ete et motivante. Je reste disponible pour un entretien.

Je vous prie d'agr\'eer, Madame, Monsieur, l'expression de mes salutations distingu\'ees.

\vspace{1cm}

\noindent\textbf{Lo\"ic Aron MBASSI EWOLO}

\end{document}
